The section passes through the midpoints of $AC,AD,BC,BD$. Its consecutive sides are parallel to $AB$ and $CD$ and have length one by the midpoint theorem. Opposite edges of a regular tetrahedron are perpendicular, so the section is a unit square of area one. Symmetry divides the original surface area equally between the pieces. Four equilateral faces of side two have total area $4\sqrt3$, so each piece receives exterior area $2\sqrt3$. Adding its new square face gives $\boxed{2\sqrt3+1}$.
\begin{asy}
texpreamble("\usepackage{amsmath,amssymb}");
size(220);
defaultpen(fontsize(9pt));
pair A=(0,3),B=(4,3),C=(0,0),D=(3,1);draw(A--B--C--A);draw(A--D--B);draw(C--D,dashed);pair P=(A+C)/2,Q=(A+D)/2,R=(B+D)/2,S=(B+C)/2;draw(P--Q--R--S--cycle,linewidth(1.2));label("$A$",A,(-1,1));label("$B$",B,(1,1));label("$C$",C,(-1,-1));label("$D$",D,(1,-1));label("$P$",P,(-1,0));label("$Q$",Q,(0,1));label("$R$",R,(1,0));label("$S$",S,(0,-1));
\end{asy}
